\setcounter{ex}{0}

% --- CÂU 1 ---
\begin{ex}
	Cho $\cos\alpha = -\dfrac{4}{5}$ với $\dfrac{\pi}{2} < \alpha < \pi$. Giá trị của $\sin 2\alpha$ bằng
	\choice
	{$\dfrac{24}{25}$}
	{\True $-\dfrac{24}{25}$}
	{$-\dfrac{7}{25}$}
	{$\dfrac{12}{25}$}
	\loigiai{
		Vì $\dfrac{\pi}{2} < \alpha < \pi$ nên $\sin\alpha > 0$. Ta có:
		$$\sin\alpha = \sqrt{1 - \cos^2\alpha} = \sqrt{1 - \left(-\dfrac{4}{5}\right)^2} = \dfrac{3}{5}.$$
		Khi đó:
		$$\sin 2\alpha = 2\sin\alpha\cos\alpha = 2 \cdot \dfrac{3}{5} \cdot \left(-\dfrac{4}{5}\right) = -\dfrac{24}{25}.$$
	}
\end{ex}

% --- CÂU 2 ---
\begin{ex}
	Cho $\tan\alpha = 3$. Giá trị của biểu thức $\dfrac{2\sin\alpha + 3\cos\alpha}{4\sin\alpha - 5\cos\alpha}$ bằng
	\choice
	{$\dfrac{8}{7}$}
	{$\dfrac{7}{9}$}
	{\True $\dfrac{9}{7}$}
	{$2$}
	\loigiai{
		Vì $\tan\alpha = 3$ nên $\cos\alpha \ne 0$. Chia cả tử và mẫu của biểu thức cho $\cos\alpha$, ta được:
		$$\dfrac{2\sin\alpha + 3\cos\alpha}{4\sin\alpha - 5\cos\alpha} = \dfrac{2\tan\alpha + 3}{4\tan\alpha - 5} = \dfrac{2 \cdot 3 + 3}{4 \cdot 3 - 5} = \dfrac{9}{7}.$$
	}
\end{ex}

% --- CÂU 3 ---
\begin{ex}
	Tập xác định của hàm số $y = \dfrac{1}{\sin x + \cos x}$ là
	\choice
	{$\mathbb{R} \setminus \left\{\dfrac{\pi}{4} + k\pi,\ k \in \mathbb{Z}\right\}$}
	{\True $\mathbb{R} \setminus \left\{-\dfrac{\pi}{4} + k\pi,\ k \in \mathbb{Z}\right\}$}
	{$\mathbb{R} \setminus \left\{-\dfrac{\pi}{4} + k2\pi,\ k \in \mathbb{Z}\right\}$}
	{$\mathbb{R} \setminus \left\{\dfrac{3\pi}{4} + k2\pi,\ k \in \mathbb{Z}\right\}$}
	\loigiai{
		Hàm số xác định khi và chỉ khi:
		$$\sin x + \cos x \ne 0 \iff \sqrt{2}\sin\left(x + \dfrac{\pi}{4}\right) \ne 0 \iff x + \dfrac{\pi}{4} \ne k\pi \iff x \ne -\dfrac{\pi}{4} + k\pi\ (k \in \mathbb{Z}).$$
		Vậy tập xác định của hàm số là $\mathscr{D} = \mathbb{R} \setminus \left\{-\dfrac{\pi}{4} + k\pi,\ k \in \mathbb{Z}\right\}$.
	}
\end{ex}

% --- CÂU 4 ---
\begin{ex}
	Hàm số nào dưới đây là hàm số chẵn?
	\choice
	{$y = x^2\sin x$}
	{$y = \sin 2x$}
	{\True $y = x\sin 2x$}
	{$y = \cos x + \sin x$}
	\loigiai{
		Xét hàm số $f(x) = x\sin 2x$ có tập xác định $\mathscr{D} = \mathbb{R}$ là tập đối xứng.\\
		Với mọi $x \in \mathbb{R}$, ta có $-x \in \mathbb{R}$ và:
		$$f(-x) = (-x)\sin(-2x) = (-x)(-\sin 2x) = x\sin 2x = f(x).$$
		Vậy $y = x\sin 2x$ là hàm số chẵn.
	}
\end{ex}

% --- CÂU 5 ---
\begin{ex}
	Nghiệm của phương trình $2\cos\left(2x - \dfrac{\pi}{6}\right) = 1$ là
	\choice
	{$x = \dfrac{\pi}{4} + k\pi$ hoặc $x = -\dfrac{\pi}{12} + k\pi\ (k \in \mathbb{Z})$}
	{$x = \dfrac{\pi}{6} + k\pi$ hoặc $x = -\dfrac{\pi}{6} + k\pi\ (k \in \mathbb{Z})$}
	{\True $x = \dfrac{\pi}{4} + k\pi$ hoặc $x = -\dfrac{\pi}{12} + k\pi\ (k \in \mathbb{Z})$}
	{$x = \dfrac{\pi}{3} + k\pi$ hoặc $x = -\dfrac{\pi}{6} + k\pi\ (k \in \mathbb{Z})$}
	\loigiai{
		Ta có:
		$$2\cos\left(2x - \dfrac{\pi}{6}\right) = 1 \iff \cos\left(2x - \dfrac{\pi}{6}\right) = \dfrac{1}{2} \iff \hoac{2x - \dfrac{\pi}{6} &= \dfrac{\pi}{3} + k2\pi \\ 2x - \dfrac{\pi}{6} &= -\dfrac{\pi}{3} + k2\pi} \iff \hoac{2x &= \dfrac{\pi}{2} + k2\pi \\ 2x &= -\dfrac{\pi}{6} + k2\pi} \iff \hoac{x &= \dfrac{\pi}{4} + k\pi \\ x &= -\dfrac{\pi}{12} + k\pi}\ (k \in \mathbb{Z}).$$
	}
\end{ex}

% --- CÂU 6 ---
\begin{ex}
	Số nghiệm của phương trình $2\sin 2x - 1 = 0$ trên đoạn $[0; 2\pi]$ là
	\choice
	{$2$}
	{$3$}
	{\True $4$}
	{$6$}
	\loigiai{
		Ta có:
		$$2\sin 2x - 1 = 0 \iff \sin 2x = \dfrac{1}{2} \iff \hoac{2x &= \dfrac{\pi}{6} + k2\pi \\ 2x &= \dfrac{5\pi}{6} + k2\pi} \iff \hoac{x &= \dfrac{\pi}{12} + k\pi \\ x &= \dfrac{5\pi}{12} + k\pi}\ (k \in \mathbb{Z}).$$
		Xét trên đoạn $[0; 2\pi]$:
		\begin{itemize}
			\item Với $x = \dfrac{\pi}{12} + k\pi \in [0; 2\pi] \iff 0 \le \dfrac{1}{12} + k \le 2 \iff -\dfrac{1}{12} \le k \le \dfrac{23}{12} \implies k \in \{0; 1\}$ ($2$ nghiệm).
			\item Với $x = \dfrac{5\pi}{12} + k\pi \in [0; 2\pi] \iff 0 \le \dfrac{5}{12} + k \le 2 \iff -\dfrac{5}{12} \le k \le \dfrac{19}{12} \implies k \in \{0; 1\}$ ($2$ nghiệm).
		\end{itemize}
		Vậy phương trình có đúng $4$ nghiệm thuộc đoạn $[0; 2\pi]$.
	}
\end{ex}

% --- CÂU 7 ---
\begin{ex}
	Giá trị lớn nhất của hàm số $y = 5\cos\left(x - \dfrac{\pi}{4}\right) + 2$ bằng
	\choice
	{$5$}
	{\True $7$}
	{$-3$}
	{$3$}
	\loigiai{
		Vì $-1 \le \cos\left(x - \dfrac{\pi}{4}\right) \le 1, \forall x \in \mathbb{R}$ nên:
		$$-3 \le 5\cos\left(x - \dfrac{\pi}{4}\right) + 2 \le 7.$$
		Vậy giá trị lớn nhất của hàm số bằng $7$ khi $\cos\left(x - \dfrac{\pi}{4}\right) = 1$.
	}
\end{ex}

% --- CÂU 8 ---
\begin{ex}
	Cho cấp số cộng $(u_n)$ có $u_3 = 7$ và $u_6 = 16$. Số hạng tổng quát của cấp số cộng là
	\choice
	{\True $u_n = 3n - 2$}
	{$u_n = 3n + 1$}
	{$u_n = 2n + 3$}
	{$u_n = 3n - 1$}
	\loigiai{
		Gọi $d$ là công sai của cấp số cộng. Theo giả thiết, ta có:
		$$\heva{u_3 &= u_1 + 2d = 7 \\ u_6 &= u_1 + 5d = 16} \iff \heva{u_1 &= 1 \\ d &= 3.}$$
		Số hạng tổng quát của cấp số cộng là:
		$$u_n = u_1 + (n - 1)d = 1 + (n - 1) \cdot 3 = 3n - 2.$$
	}
\end{ex}

% --- CÂU 9 ---
\begin{ex}
	Cho cấp số cộng $(u_n)$ có $u_1 = -2$ và công sai $d = 3$. Tổng $15$ số hạng đầu bằng
	\choice
	{$275$}
	{\True $285$}
	{$295$}
	{$315$}
	\loigiai{
		Tổng $15$ số hạng đầu của cấp số cộng là:
		$$S_{15} = 15u_1 + \dfrac{15 \cdot 14}{2}d = 15 \cdot (-2) + 105 \cdot 3 = -30 + 315 = 285.$$
	}
\end{ex}

% --- CÂU 10 ---
\begin{ex}
	Cho hình chóp $S.ABCD$ có đáy $ABCD$ là hình bình hành tâm $O$, $N$ là trung điểm của $SD$ (hình vẽ dùng chung cho câu 10, 11, 12). Giao tuyến của hai mặt phẳng $(SAB)$ và $(SCD)$ là đường thẳng đi qua $S$ và
	\begin{center}
		\begin{tikzpicture}[scale=0.85, font=\footnotesize, line join=round, line cap=round]
			\coordinate (A) at (0,0);
			\coordinate (B) at (3.2,0);
			\coordinate (D) at (1.3,1.1);
			\coordinate (C) at (4.5,1.1);
			\coordinate (O) at ($(A)!0.5!(C)$);
			\coordinate (S) at (2.2,3.4);
			\coordinate (N) at ($(S)!0.5!(D)$);
			\draw[thick] (S)--(A) (S)--(B) (S)--(C) (A)--(B)--(C);
			\draw[dashed] (S)--(D) (A)--(D)--(C) (A)--(C) (B)--(D);
			\foreach \p/\pos in {S/above, A/below left, B/below, C/right, D/above left, O/below, N/above left}
				\node[\pos] at (\p) {$\p$};
			\foreach \p in {S,A,B,C,D,O,N}
				\fill (\p) circle (1.2pt);
		\end{tikzpicture}
	\end{center}
	\choice
	{song song với $AD$}
	{\True song song với $AB$}
	{song song với $BD$}
	{đi qua điểm $O$}
	\loigiai{
		Hai mặt phẳng $(SAB)$ và $(SCD)$ có điểm chung $S$ và lần lượt chứa hai đường thẳng song song $AB \parallel CD$.\\
		Do đó giao tuyến của $(SAB)$ và $(SCD)$ là đường thẳng qua $S$ song song với $AB$ (và $CD$).
	}
\end{ex}

% --- CÂU 11 ---
\begin{ex}
	Gọi $K$ là giao điểm của đường thẳng $SC$ và mặt phẳng $(ABN)$. Khẳng định nào sau đây đúng?
	\choice
	{$KC = 2SK$}
	{\True $K$ là trung điểm của $SC$}
	{$SK = 2KC$}
	{$SK = 3KC$}
	\loigiai{
		Hai mặt phẳng $(ABN)$ và $(SCD)$ có điểm chung $N$ và lần lượt chứa hai đường thẳng song song $AB \parallel CD$.\\
		Suy ra giao tuyến của $(ABN)$ và $(SCD)$ là đường thẳng qua $N$ song song với $CD$, cắt $SC$ tại $K$.\\
		Trong tam giác $SCD$, vì $N$ là trung điểm của $SD$ và $NK \parallel CD\ (K \in SC)$ nên $K$ là trung điểm của $SC$.
	}
\end{ex}

% --- CÂU 12 ---
\begin{ex}
	Khẳng định nào sau đây đúng?
	\choice
	{$ON \parallel (SBC)$}
	{$ON \parallel (ABCD)$}
	{\True $ON \parallel (SAB)$}
	{$ON \parallel (SCD)$}
	\loigiai{
		Xét tam giác $SBD$ có $O$ là trung điểm của $BD$ và $N$ là trung điểm của $SD$.\\
		Suy ra $ON$ là đường trung bình của tam giác $SBD \implies ON \parallel SB$.\\
		Mà $SB \subset (SAB)$ và $ON \not\subset (SAB)$ nên $ON \parallel (SAB)$.
	}
\end{ex}
